\section{User Feedback Gathering}
\label{research:UserFeedbackGathering}
\textbf{Source:} \cite{HowWriteFeedback}

Feedback forms are valuable tools used by companies to gather information about their performance, products, and services. They help organizations make data-driven decisions by collecting insights from customers, employees, or other stakeholders. When used effectively, feedback forms highlight strengths and weaknesses, and encourage continuous improvement. By understanding how to design and use these forms, individuals can contribute meaningful data that informs strategic planning and professional growth.

\begin{enumerate}
    \item It’s important to clearly define the purpose of the form and tailor questions to the target audience. Using a variety of question types, such as open-ended responses and rating scales, can provide both detailed insights and easily measurable data. 
    \item Good formatting, consistency, and unbiased wording ensure the form is professional and encourages honest, thoughtful responses. 
    \item Organizing questions into categories and collecting relevant demographic information can also improve data analysis and interpretation.
    \item proofreading and testing the form before distribution ensures clarity and usability. Depending on the goal, feedback forms can be anonymous to encourage honesty
    \item Digital forms offer advantages in scalability and data collection especially in software products such as mine
\end{enumerate}
